% ============================================================
% Section 121: 一维晶体能带与有效质量计算
% ============================================================
\section{固体物理：一维晶体能带参数计算}
\label{sec:1d-band-params}

\begin{modelbox}[题目：一维晶体的禁带宽度与有效质量]
设晶格常数为 $a$ 的一维晶体，导带极小值附近能量 $E_C(k)$ 为
\[
E_C(k)=\frac{\hbar^2k^2}{3m_0}+\frac{\hbar^2(k-k_1)^2}{m_0},
\]
价带极大值附近能量 $E_V(k)$ 为
\[
E_V(k)=\frac{\hbar^2k_1^2}{6m_0}-\frac{3\hbar^2k^2}{m_0},
\]
其中 $m_0$ 为电子质量，$k_1=\dfrac{1}{2a}$，$a=3.14\times10^{-10}\,\text{m}$。试求：

(1) 禁带宽度；

(2) 导带底电子有效质量；

(3) 价带顶电子有效质量；

(4) 价带顶电子跃迁到导带底的准动量的变化。
\end{modelbox}

\begin{tipbox}[考点快速分析]
\begin{itemize}
    \item \textbf{能带极值位置}：由 $\dd{E}{k}=0$ 确定
    \item \textbf{禁带宽度}：$E_g=E_{C,\min}-E_{V,\max}$
    \item \textbf{有效质量}：$m^*=\hbar^2\left(\dd[2]{E}{k}\right)^{-1}$，导带底为正，价带顶为负
    \item \textbf{准动量变化}：$\Delta p=\hbar\Delta k$，注意直接/间接跃迁的区别
\end{itemize}
\end{tipbox}

\begin{derivationbox}[标准解题过程]
\textbf{(1) 禁带宽度 $E_g$}

禁带宽度定义为导带底能量与价带顶能量之差：$E_g=E_{C,\min}-E_{V,\max}$。

\textit{导带底位置}：
\[
\dd{E_C}{k}=\frac{\hbar^2}{m_0}\left(\frac{2k}{3}+2(k-k_1)\right)=0
\implies \frac{8k}{3}=2k_1
\implies k_c=\frac{3}{4}k_1.
\]
代入得导带底能量：
\[
E_{C,\min}=\frac{\hbar^2}{m_0}\left[\frac{1}{3}\left(\frac{3k_1}{4}\right)^2+\left(\frac{3k_1}{4}-k_1\right)^2\right]
=\frac{\hbar^2k_1^2}{m_0}\left(\frac{3}{16}+\frac{1}{16}\right)=\frac{\hbar^2k_1^2}{4m_0}.
\]

\textit{价带顶位置}：
$E_V(k)$ 为开口向下的抛物线，极大值在 $k_v=0$ 处：
\[
E_{V,\max}=\frac{\hbar^2k_1^2}{6m_0}.
\]

因此禁带宽度：
\[
E_g=\frac{\hbar^2k_1^2}{4m_0}-\frac{\hbar^2k_1^2}{6m_0}=\frac{\hbar^2k_1^2}{12m_0}=\frac{h^2}{48m_0a^2}.
\]
代入数值 $h=6.63\times10^{-34}\,\text{J\cdot s}$，$m_0=9.1\times10^{-31}\,\text{kg}$，$a=3.14\times10^{-10}\,\text{m}$：
\begin{equation}
\boxed{E_g\approx1.02\times10^{-19}\,\text{J}\approx0.64\,\text{eV}.}
\end{equation}

\textbf{(2) 导带底电子有效质量 $m_c^*$}

对 $E_C(k)$ 求二阶导数：
\[
\dd[2]{E_C}{k}=\frac{\hbar^2}{m_0}\left(\frac{2}{3}+2\right)=\frac{8\hbar^2}{3m_0}.
\]
因此
\begin{equation}
\boxed{m_c^*=\hbar^2\left(\frac{8\hbar^2}{3m_0}\right)^{-1}=\frac{3}{8}m_0=0.375\,m_0.}
\end{equation}

\textbf{(3) 价带顶电子有效质量 $m_v^*$}

对 $E_V(k)$ 求二阶导数：
\[
\dd[2]{E_V}{k}=-\frac{6\hbar^2}{m_0}.
\]
因此
\begin{equation}
\boxed{m_v^*=\hbar^2\left(-\frac{6\hbar^2}{m_0}\right)^{-1}=-\frac{1}{6}m_0\approx-0.167\,m_0.}
\end{equation}
价带顶电子有效质量为负值，表明其行为与正电荷类似。空穴的有效质量为 $m_h^*=-m_v^*=\dfrac{1}{6}m_0$。

\textbf{(4) 跃迁时准动量的变化}

准动量 $p=\hbar k$。价带顶电子波矢 $k_v=0$，导带底电子波矢 $k_c=\dfrac{3}{4}k_1$。跃迁时准动量的变化：
\[
\Delta p=\hbar(k_c-k_v)=\frac{3}{4}\hbar k_1=\frac{3}{4}\cdot\frac{h}{2\pi}\cdot\frac{1}{2a}=\frac{3h}{16\pi a}.
\]
代入数值：
\begin{equation}
\boxed{\Delta p\approx7.9\times10^{-25}\,\text{kg\cdot m/s}.}
\end{equation}
\end{derivationbox}

\begin{formulabox}[答案汇总]
\begin{center}
\begin{tabular}{ll}
\toprule
\textbf{问题} & \textbf{答案} \\
\midrule
(1) 禁带宽度 $E_g$ & $0.64\,\text{eV}$ \\
(2) 导带底电子有效质量 $m_c^*$ & $\dfrac{3}{8}m_0=0.375\,m_0$ \\
(3) 价带顶电子有效质量 $m_v^*$ & $-\dfrac{1}{6}m_0\approx-0.167\,m_0$ \\
(4) 准动量变化 $\Delta p$ & $7.9\times10^{-25}\,\text{kg\cdot m/s}$ \\
\bottomrule
\end{tabular}
\end{center}
\end{formulabox}

\begin{insightbox}[物理图像]
\begin{itemize}
    \item 一维能带的极值位置由能量对波矢的一阶导数为零确定，这是求有效质量和禁带宽度的起点。
    \item 导带底曲率为正，有效质量为正；价带顶曲率为负，有效质量为负。负有效质量意味着电子在价带顶受到的外场加速度与力方向相反，等价于带正电的空穴。
    \item 本题中导带底与价带顶在 $k$ 空间不重合（$k_c\neq k_v$），对应间接带隙。直接的光学跃迁需满足 $\Delta k=0$，因此间接跃迁需要声子或缺陷辅助以补偿动量差。
\end{itemize}
\end{insightbox}
